\chapter*{ПРИЛОЖЕНИЕ В}
\addcontentsline{toc}{chapter}{ПРИЛОЖЕНИЕ В}
\begin{lstlisting}[caption={Определение множества всех слов в наборе рецептов и количеств каждого слова}, label={lst:rvec_1}]
import pandas as pd
import ast

def clean(s):
    result = ''
    s = s.lower()
    for c in s:
        if c >= 'a' and c <= 'z' or c >= '0' and c <= '9':
            result += c
    return str(result)

def words(s):
    i = 0
    result = []
    while i < len(s):
        if not s[i].isalpha():
            i += 1
        else:
            j =  i + 1
            while j < len(s) and s[j].isalpha():
                j += 1
            w = clean(s[i:j])
            if len(w) > 0:
                result.append(w)
            i = j + 1
    return result

def get_document(raw_recipes, i):
    name = raw_recipes['name'][i]
    if type(name) == float:
        name = ''
    description = raw_recipes['description'][i]
    if type(description) == float:
        description = ''    
    steps = str(raw_recipes['steps'][i])
    tags = ast.literal_eval(raw_recipes['tags'][i])
    ingredients = ast.literal_eval(raw_recipes['ingredients'][i])
    all_words = words(name)
    all_words += words(description)
    for elem in tags:
        all_words += words(elem)
    for elem in ingredients:
        all_words += words(elem)
    all_words += words(steps)
    return all_words

BASE_PATH = """/kaggle/input/datasets/shuyangli94/
food-com-recipes-and-user-interactions/"""
raw_recipes = pd.read_csv(BASE_PATH + 'RAW_recipes.csv')
n = len(list(raw_recipes['id']))
documents = []
for i in range(n):
    documents.append(get_document(raw_recipes, i))

word_frequency = dict()
for doc in documents:
    for elem in doc:
        if elem in word_frequency:
            word_frequency[elem] += 1
        else:
            word_frequency[elem] = 1
count_word_pairs = sorted([(b, a) for a, b in word_frequency.items()])
\end{lstlisting}

\begin{lstlisting}[caption={Определение набора рассматриваемых слов}, label={lst:rvec_2}]
candidate_words = dict([(b, a) for a, b in count_word_pairs[85000 : 85000 + 4120]])
stop_words = {
    # Грамматический мусор
    'the', 'and', 'with', 'for', 'from', 'you', 'off', 'this', 'that', 'your',
    'but', 'not', 'are', 'was', 'were', 'have', 'has', 'had', 'can', 'will',
    'about', 'all', 'out', 'into', 'then', 'them', 'they', 'their', 'than',
    
    # Общие глаголы готовки (они есть везде и не отражают вкус)
    'cook', 'serve', 'drink', 'make', 'made', 'add', 'added', 'adding', 'use', 
    'used', 'using', 'take', 'put', 'place', 'bring', 'keep', 'let', 'get', 'comes',
    
    # Меры веса и времени (шум для модели предпочтений)
    'minutes', 'minute', 'hours', 'hour', 'grams', 'gram', 'cups', 'cup', 
    'tbsp', 'teaspoon', 'tsp', 'ounces', 'ounce', 'inch', 'inches', 'size',
    'medium', 'small', 'large', 'pieces', 'piece', 'amount', 'weight',
    
    # Оценочные слова (люди пишут их в комментариях)
    'unfortunately', 'favourite', 'favorite', 'meal', 'food', 'tasty', 
    'calorie', 'great', 'like', 'both', 'good', 'best', 'delicious', 'love',
    'easy', 'perfect', 'wonderful', 'nice', 'recipe', 'recipes', 'dish', 'friendly'
    
    # Другой текстовый шум
    'doesnt', 'appeared', 'beyond', 'forth', 'unexpected', 'more', 'big',
    'inspiration', 'previous', 'possibly', 'watcher', 'modifications', 'low', 'less', 'until',
    'one', 'two', 'three', 'four', 'five', 'six', 'seven', 'eight', 'nine', 'ten',
    'many', 'few', 'pour', 'spoon', 'fork', 'way', 'brother', 'sister', 'father', 'mother',
    'supposed', 'none', 'pkg', 'boy', 'girl', 'naan'
}
candidate_words = dict([(b, a) for a, b in count_word_pairs[86000:]])
selected_words = dict()
for w, count in candidate_words.items():
    # Убираем слишком короткие слова, стоп-слова и проверяем, что это не число
    if len(w) > 2 and (w not in stop_words) and not w.isdigit():
        selected_words[w] = count

# Сортируем от самых частых к менее частым и берем топ-3500
final_vocabulary = dict(sorted(selected_words.items(), key=lambda item: item[1], reverse=True)[:3500])
final_vocabulary # Набор рассматриваемых слов.
\end{lstlisting}

\begin{lstlisting}[caption={Преобразование слов в числовые токены для ускорения работы}, label={lst:rvec_3}]
ALPHABET = "abcdefghijklmnopqrstuvwxyz0123456789"
ALPH_LEN = len(ALPHABET)
letter_hash = dict()
for i, c in enumerate(ALPHABET):
    letter_hash[c] = i + 1

def token(s):
    M = 2_147_483_647
    result = 0
    for c in s:
        result = (result * ALPH_LEN + letter_hash[c]) % M
    return result

token_lists = [[token(w) for w in documents[i]] for i in range(len(documents))]
token_sets = [set(doc) for doc in token_lists]
vacabulary_tokens = [token(w) for w in final_vocabulary]
\end{lstlisting}

\begin{lstlisting}[caption={Получение векторов рецептов}, label={lst:rvec_4}]
def doc_count(w):
    result = 0
    for token_set in token_sets:
        if w in token_set:
            result += 1
    return result

document_count_of = dict()
for i, t in enumerate(vacabulary_tokens):
    document_count_of[t] = doc_count(t)

def tf_idf(w, doc):
    tf = doc.count(w) / (len(doc) + 0.001)
    idf = np.log(len(documents) / (document_count_of[w] + 0.001))
    return tf * idf

def documents_to_recipe_vec(documents):
    n, m = len(documents), len(vacabulary_tokens)
    recipe_vec = []
    for i in range(n):
        recipe_vec.append([])
        for j in range(m):
            vec = tf_idf(vacabulary_tokens[j], documents[i])
            recipe_vec[i].append(vec)
        if i % 1000 == 0:
            print("documents_to_recipe_vec: processed recipes =", i + 1)
    return np.array(recipe_vec)
 
 recipe_vectors = documents_to_recipe_vec(token_lists)
 recipe_vectors # Результирующие вектора всех рецептов.
\end{lstlisting}
